% CLAUDE REPORT::WORSTCASE-DERIV-0001::2026-11-06
% !TEX program = xelatex
\documentclass[11pt]{article}
\usepackage[margin=1in]{geometry}
\usepackage{amsmath,amssymb,amsthm}
\usepackage{xurl}
\newtheorem{theorem}{Theorem}
\newtheorem{proposition}[theorem]{Proposition}
\newtheorem{corollary}[theorem]{Corollary}
\theoremstyle{remark}\newtheorem*{remark}{Remark}
\newcommand{\status}[1]{\textbf{[#1]}}
\newcommand{\UNSAT}{\mathrm{UNSAT}}
\title{CLAUDE WORSTCASE-DERIV-0001\\[2pt]\large The worst-case minimum-derivation functional $\overline W_R$:\\ what the bridge $\mathrm{SAT}\in P\Rightarrow \overline W_R\le\mathrm{poly}$ can and cannot be}
\author{CLAUDE (scheduled run; additive, not an LSPV run)}
\date{2026-11-06}
\begin{document}\maketitle
\noindent\textbf{Status.} P versus NP remains OPEN. This report proves no lower bound. It classifies the proposed target
$\overline W_R(N)=\max_{F\in\UNSAT,|F|\le N}\min_{\pi:F\vdash_R\bot}\mathrm{cost}_R(\pi)$ and the proposed bridge, using standard
theorems (all \status{KNOWN THEOREM} unless marked). No schedules, allocations, ownership or root \texttt{main.tex} changed.

\section{Free strategy discovery is legitimate, and costs nothing here}
Taking the min over all derivations without charging search is exactly the Cook--Reckhow measure: $R$ is a proof system and
$\overline W_R$ is its optimal proof size. \status{KNOWN THEOREM} (Cook--Reckhow 1979): $\overline W_R\le\mathrm{poly}$ for a sound, complete,
poly-time-verifiable $R$ iff $R$ is polynomially bounded, and a polynomially bounded $R$ exists iff $\mathrm{NP}=\mathrm{coNP}$.
Hence \emph{every} superpolynomial lower bound for a fixed $R$ is consistent with $\mathrm{P}=\mathrm{NP}$ unless a bridge is proved.

\section{Three bridge candidates, ranked}
\begin{proposition}[Unconditional strong form is equivalent to NP$\ne$coNP]\status{KNOWN THEOREM}
``$\overline W_R$ superpolynomial for \emph{all} sound $R$'' $\iff$ $\mathrm{NP}\neq\mathrm{coNP}$. This is strictly stronger than $\mathrm{P}\neq\mathrm{NP}$,
so the all-$R$ route proves more than needed and is not the weakest sufficient statement.
\end{proposition}
\begin{proposition}[Bridge for fixed natural $R$ is itself open]\status{OPEN}
For $R\in\{$Frege, Extended Frege, $\ldots\}$, ``$\mathrm{SAT}\in P\Rightarrow R$ polynomially bounded'' is not known: $\mathrm{P}=\mathrm{NP}$ gives a poly-time
decider $M$ but $R$ need not be able to certify $M$'s correctness. Hence a fixed-$R$ bridge cannot be assumed as a lemma.
\end{proposition}
\begin{theorem}[Provable-correctness bridge]\status{KNOWN THEOREM} (Krajicek--Pudlak style; the standard route)
Let $T$ be a fixed sound theory and $R_T$ the proof system ``$F$ is unsatisfiable iff $T$ proves, via the propositional translation, $\neg F$''.
If $T\vdash$ ``$M$ decides SAT in time $n^c$'' for some fixed $M,c$, then $R_T$ is polynomially bounded.
\end{theorem}
\emph{Sketch.} The constant-size proof of correctness of $M$ is reused for every $F$; the length-$\mathrm{poly}(|F|)$ computation trace of $M(F)$ plus
reflection gives a $T$-proof of $\neg F$ of size $\mathrm{poly}(|F|)$.\hfill$\square$

\begin{corollary}[Exact disjunction, not P$\ne$NP]
If one proves $\overline W_{R_T}$ superpolynomial, then: \emph{either} $\mathrm{P}\ne\mathrm{NP}$ \emph{or} $\mathrm{P}=\mathrm{NP}$ holds but no algorithm for it is
provably correct in $T$. Closing the second disjunct is the missing load-bearing arrow. It is a ``no independence from $T$ of the correctness of a poly-time SAT algorithm''
statement; it is not implied by the lower bound and is the precise form of the weakest bridge.
\end{corollary}
\begin{remark} This answers the agents' ``weakest sufficient bridge'': for a \emph{fixed} $R$ no unconditional bridge is available; the usable one is
$R_T$ plus the hypothesis ``some poly-time SAT algorithm is $T$-provably correct'' (Levin-optimal-search flavour). Making it unconditional requires either
an optimal proof system (open) or a proof that independence of this kind is impossible (open; cf. Krajicek--Pudlak optimal-proof-system conjectures).
\end{remark}

\section{Role of the five agents under this analysis}
\begin{itemize}
\item \textbf{Lower-bound lead:} target $\overline W_{R}$ for the strongest $R$ with known lower-bound technology; the load-bearing question is which $R$ admits Theorem 3's bridge with $T$ weak enough to be attackable.
\item \textbf{Circuit-reduction builder:} a fixed combined system containing Res, cutting planes/algebraic rules and circuit compression is only useful if it is still verifiable in poly time (else $\mathrm{cost}_R$ not a proof system). Strengthening $R$ helps the bridge and hurts the lower bound.
\item \textbf{Hypergraph/certificates:} hard UNSAT families must beat Res, PC/Nullstellensatz, cutting planes, bounded-depth Frege; all currently proven bounds stop below Frege (AC$^0[p]$-Frege lower bounds still open).
\item \textbf{Auditor:} cost must be size of a poly-time-checkable object; peak memory/time of \emph{executing} the DAG is $\ge$ DAG size, so time lower bounds from $\overline W_R$ transfer only for algorithms \emph{that output derivations}, not arbitrary deciders.
\item \textbf{Red team:} the decisive attack is the independence disjunct above and the Cook--Reckhow observation that fixed-$R$ bounds are consistent with $\mathrm{P}=\mathrm{NP}$.
\end{itemize}

\section{Literature check (this run)}
Recent primary items found (not verified line-by-line here; \status{CANDIDATE} relevance only): Davis--Robere, Res(log) proves bounded-depth Frege lower bounds (ECCC TR26-055);
semantic tree-like Frege bounded-line-size lower bounds (arXiv 2604.28172); Bhattacharya--Chattopadhyay, Res($\oplus$) (STOC 2026, arXiv 2507.23008);
CCC 2026 paper 8 on prover--adversary games and a conditional Extended Frege bound. None gives a superpolynomial Frege bound or touches the bridge.
Classification: useful analogies (self-proving-lower-bound systems resemble the reflection in Theorem 3), no transferable theorem.

\section{Next target (one)}
State and prove a quantitative version of Theorem 3 with explicit overhead for the repository's own system (signed multi-sign hypergraph DAG inference), identifying the weakest $T$; then ask whether any known hard family is hard for $R_T$.

\paragraph{Honest summary.} No proof of $\mathrm{P}\ne\mathrm{NP}$ is claimed. Load-bearing open arrows: (i) superpolynomial $\overline W_R$ for a system strong enough for the bridge; (ii) the bridge itself without provable-correctness hypothesis.
% END_OF_VERIFIED_REPORT::CLAUDE::WORSTCASE-DERIV-0001
\end{document}
